\chapter*{Travaux Pratiques - Les boucles }

\begin{Exercise}[title={Echauffement}]
Créez un nouveau fichier python que vous nommerez \textit{echauffement}.\\
Reprenons l'exercice 3 de la feuille d'exercice où vous deviez afficher un sapin de noël comme ci-dessous :
\begin{center}
    *\\
    *\\
   ***\\
  *****\\
 *******\\
   ***\\
   ***\\
\end{center}
\Question Sur le même principe faites le motif suivant :
\begin{center}
******\\
 ****\\
  **\\
  **\\
 ****\\
******\\
\end{center}
\end{Exercise}

\begin{Exercise}[title={Le juste prix}]
Créez un nouveau fichier python que vous nommerez \textit{LeJustePrix}.\\
Nous allons créer le jeu du juste prix. Le programme va choisir au hasard un prix entre 1€ et 100€ et vous devez retrouver ce prix. Vous avez 10 essais. A chaque essai le programme vous dit si vous êtes trop haut ou trop bas ou si vous avez gagné. Si au bout de 10 essais vous n'avez pas trouvé le prix, vous avez perdu. \\
Répondez aux questions qui vous guident pour créer ce jeu.
\Question Importer le module random puis enregistrer dans une variable le prix tiré au hasard entre 1 et 10. Pour tirez au hasard, vous utilisez la fonction \texttt{randint(nombre\_min,nombre\_max)}. Remplcez le nombre min et max par les valeurs données dans l'énoncé.
\Question Utilisez une boucle qui va exécuter les 10 essais. Dans ce bloc, vous demandez à l'utilisateur de faire une proposition de nombre. Tester votre programme.
\Question Ajouter les conditions qui disent à l'utilisateur si le nombre saisi est trop haut, trop bas ou si c'est la bonne réponse. Tester votre programme.
\Question Changez le \texttt{randint(1,10)} en \texttt{randint(1,100)}. Tester votre programme.
\Question Ajouter la ligne du code qui dit à l'utilisateur qu'il a perdu s'il n'a pas trouvé le bon chiffre au bout des 10 essais. Tester votre programme.
\Question Si besoin changer votre boucle \texttt{for} en une boucle \texttt{while}, pour pouvoir faire le cas où l'utilisateur trouve le bon chiffre avant la fin des 10 essais. Tester votre programme.
    
\end{Exercise}

\begin{Exercise}[title={Dés bloqués remix}]
Créez un nouveau fichier python que vous nommerez \textit{DesBloques}.\\
Ce jeu se joue avec 5 dés. Le but est de marquer le plus de points possibles en 10 tours. Voici les règles. Vous lancez les 5 dés.
\begin{itemize}
    \item Si vous avez obtenu deux 1, vous ne marquez aucun points pour ce lancer.
    \item Si vous n'avez pas obtenu deux 1, additionnez le total des dés (1 exclu) et mémorisez-le.
\end{itemize}
Vous pouvez relancer le nombre de dés que vous voulez, en laissant le 1 de côté et recalculer votre points à chaque lancer mais attention dès que deux de vos dés sont à 1, vous perdez tout ! A vous de vous arrêtez au bon moment.\\

%
\textbf{Déroulement d'une manche : on lance des dés tant qu'on a pas deux dés 1 ou qu'on ne choisit pas de s'arrêter}
\Question Créez 5 variables qui vont enregistre la valeur des dés et simuler un premier lancer avec la fonction \texttt{randint(1,6)} comme vu pendant les exercices sur les conditions. Testez (=exécutez) votre programme.
\Question Affichez les valeurs des dés. Tester votre programme.
\Question Comparez si deux de vos dés ou plus sont des 1. Si oui, mettez le score à 0 et afficher un message avec le score. Tester votre programme.
\Question Créer une variable booléenne \texttt{stop}. Affecter lui la valeur True si deux de vos dés ou plus sont des 1. Tester votre programme.
\Question Sinon, faite la somme des dés (1 exclu) et demandez à l'utilisateur pour chaque dé (sauf les 1) s'il veut le relancer. Tester votre programme.
\Question Relancer les dés et refaite la somme (attention les scores ne s'additionnent pas, on garde le dernier lancer à chaque fois). Tester votre programme.
\Question Passer la variable booléenne à True dès qu'il y a deux dés à 1.
\Question Si l'utilisateur ne veut relancer aucun dés, affecter True à la variable stop.
\Question Re-exécutez cette partie tant que stop n'est pas True (utilisation d'une boucle while). Tester votre programme.\\

\noindent\textbf{Déroulement de 3 manches, on additionne le score de chaque manche}
\Question Faites tout ceci 3 fois à l'aide d'une boucle for, et additionnez les scores à chaque fois pour le stocker dans le score total.
\Question Affichez le score total à la fin de chaque manche.

\Question Ajouter un joueur.
\Question Affichez le vainqueur après les manches. Tester pour 3 manches.
\Question Testez pour 10 manches.
    
\end{Exercise}